% Georgian abstract — machine-assisted translation of the English abstract.
% VERIFY/POLISH with a native speaker before final submission.
\begingroup\georgianfont
მოწყობილობაზე მომუშავე ასისტენტებმა მომხმარებლის სახელით უნდა შეასრულონ
მოქმედებები --- დააყენონ მაღვიძარა, შეცვალონ პარამეტრი, დარეკონ --- ისე, რომ
მოთხოვნა ღრუბელში არ გააგზავნონ. ეს ნაშრომი იკვლევს, თუ რამდენად დიდი ენობრივი
მოდელია რეალურად საჭირო ასეთი \emph{მოქმედებისთვის}. ჩვენ წარმოგიდგენთ
A.L.F.R.E.D.-ს --- ლოკალურ, ორდონიან შუამავალ სისტემას, რომელიც თითოეულ
მოთხოვნას ანაწილებს პატარა, მოწყობილობაზე მომუშავე ``რეფლექსორსა'' (2B მოდელი,
რომელიც ერთი გავლით ავსებს დისტილირებულ ბრძანების შაბლონს) და უფრო მძიმე
``მოაზროვნეს'' (აგენტური ციკლი) შორის. ოთხი შედეგი ადასტურებს მთავარ დებულებას.
პირველი: ბრძანების აგებისას 2B რეფლექსორი შაბლონით უტოლდება 17-ჯერ უფრო დიდ
თავისუფალ მოდელს (ორივესთვის 100\%), ხოლო იგივე 2B შაბლონის გარეშე აღწევს მხოლოდ
30\%-ს --- გადამწყვეტია დისტილაცია და არა მასშტაბი. მეორე: აბსტრაქციის სარგებელი
იცვლება ინტერფეისის სირთულის მიხედვით --- სუფთა API-ზე ეს ეფექტურობის მოგებაა,
ხოლო რეალისტურ API-ზე --- სიზუსტის. მესამე და მთავარი: 35B მოდელიც კი ვერ ხვდება
მოწყობილობის ``ამოუცნობ'' კონვენციებს; მასშტაბი იძლევა აღმოჩენას (რომელი გასაღები),
მაგრამ არა კონვენციას (რას ნიშნავს მნიშვნელობა), რომელსაც მხოლოდ დისტილირებული
შაბლონი ატარებს. მეოთხე: ასეთი შაბლონი შეიძლება ავტომატურად შექმნას ``მასწავლებელმა''
მოდელმა, რომელიც დავალებებს არ ხედავს, რაც 2B რეფლექსორს 25\%-დან 76--80\%-მდე
ზრდის; ამასთან, მასწავლებლის ხარისხი პირდაპირ აისახება მოსწავლეზე (ღრუბლოვანი 80\%
ლოკალური 60\%-ის წინააღმდეგ). დარჩენილი ღია შემზღუდველი არის მარშრუტიზაცია და არა
შესრულება. ნაშრომი ამ შედეგებს აყალიბებს როგორც ``შემსრულებლის პარადიგმას'': ძლიერი
მასწავლებლით ერთხელ დავადისტილიროთ მოწყობილობის კონვენციები და შემდეგ იაფად,
სამუდამოდ გამოვიყენოთ პატარა, მოწყობილობაზე მომუშავე რეფლექსორით.

\vspace{0.6em}
\noindent\textbf{საკვანძო სიტყვები:} მოწყობილობაზე მომუშავე ენობრივი მოდელები,
ცოდნის დისტილაცია, აგენტების მარშრუტიზაცია, ხელსაწყოების გამოყენება, განზრახვის
შესრულება, კონფიდენციალურობა.
\endgroup
